% ===================================================================
%  اللوحة نمرة ١٥ — السوق. --- المأكولات.
%
%  Traduction, faite en 2026, en arabe egyptien ecrit, sur l'ido de
%  texto/io. Regles de la colonne : voir 10-lawha-01.tex. Une seule
%  numerotation.
%
%  C'EST LE TABLEAU OU LES DEUX ARABES SE SEPARENT LE PLUS, ET SUR CE
%  QUI SE MANGE — exactement comme le cantonais et le mandarin au
%  meme tableau.
%
%    توم (101)    l'ail        MSA : ثوم   — la regle ث → ت, encore
%    رز (58)      le riz       MSA : أرز
%    بطاطس (61)   les pommes de terre  MSA : بطاطا
%    قرنبيط (15)  le chou-fleur
%    كرنب (17)    le chou
%    كوسة (19)    la courgette
%    بسلّة         les petits pois
%    جرجير        le cresson
%    كرّات         les poireaux
%    بقدونس       le persil
%    خرشوف (23)   l'artichaut
%    قراصيا (65)  les pruneaux
%
%  ET LA MER DONNE SES NOMS PROPRES :
%
%    سمك موسى     la sole — « le poisson de Moise »
%    حنكليس       l'anguille
%    بلح البحر    les moules — « les dattes de mer »
%    جندوفلي (98) les huitres
%    كابوريا      le crabe
%    استاكوزا     la langouste
%    رنجة (60)    le hareng fume — celui de Cham el-Nessim
%    سمّان (42)    les cailles
%
%  Trois de ces noms sont des images devenues des noms : le poisson
%  de Moise, les dattes de mer, et le عيش الغراب du tableau 9. C'est
%  le meme procede que عين الجمل et أبو فروة — l'egyptien nomme par
%  ressemblance et cesse ensuite d'y penser.
%
%  LE PAIN REVIENT SOUS SA FORME LA PLUS EGYPTIENNE : العيش الفينو
%  (82), le petit pain blanc, a cote du عيش ordinaire du tableau 6.
%  Et le boulanger est الفرّان, l'epicier البقّال, le patissier
%  الحلواني.
%
%  DEUX METIERS DE PLUS EN -اتي ET -جي : المسنّاتي (31) le
%  remouleur, المكوجيّة (86) la repasseuse. Cela porte a onze le
%  compte des noms de metier ainsi formes depuis le tableau 4.
%
%  LE TABLEAU N'A DEMANDE AUCUNE REPRISE. Il enumere, comme le
%  tableau 15 cantonais, et pour la meme raison : une enumeration n'a
%  pas de possesseur a placer.
% ===================================================================

%%K t15-apar-1 apar
\VUsaut{4.0mm}
\VUtitre{15.0pt}{60}{اللوحة نمرة ١٥}
\VUsaut{3.0mm}

%%K t15-tit-1 sub
\VUpk{11.4pt}{40}{السوق. --- المأكولات.}
\VUsaut{3.0mm}

%%K t15-01-1 p
1. --- أغلب التجّار اللي بيوصّلوا لنا البضايع اللازمة لأكلنا بيتجمّعوا تحت
\VUgras{شادر} \textsuperscript{(1)} \VUgras{السوق المغطّى} أو في الشوارع اللي جنبه.

%%K t15-02-1 p
2. --- و\VUgras{بتاع لحمة الخنزير} \textsuperscript{(2)}، اللي بيبيع
\VUgras{جامبون} \textsuperscript{(3)}، و\VUgras{سجق أسود} \textsuperscript{(4)}،
و\VUgras{سجق} \textsuperscript{(5)}، و\VUgras{ممبار} \textsuperscript{(6)}،
و\VUgras{شحمة} \textsuperscript{(7)}، واقف جنب \VUgras{بتاعة الزبدة والجبنة}
\textsuperscript{(8)}، اللي فرشتها مليانة \VUgras{جبنة هولاندي} \textsuperscript{(9)}
بقشرة حمرا، و\VUgras{جبنة كامامبير} \textsuperscript{(10)}، و\VUgras{أقراص جبنة
جرويير} \textsuperscript{(11)}، و\VUgras{قوالب زبدة} \textsuperscript{(12)} طازة.

%%K t15-03-1 p
3. --- وقدامها، \VUgras{بتاعة الخضار} \textsuperscript{(13)} بتعرض على زبوناتها
\VUgras{طماطم} \textsuperscript{(14)} حلوة، و\VUgras{قرنبيط} \textsuperscript{(15)}،
و\VUgras{خس} \textsuperscript{(16)}، و\VUgras{كرنب} \textsuperscript{(17)}،
و\VUgras{كوسة} \textsuperscript{(19)}، و\VUgras{كانتالوب} \textsuperscript{(18)}، وحتى
\VUgras{عيش غراب} \textsuperscript{(20)}. بس \VUgras{بتاع البيرة}، اللي وقّف
\VUgras{عربية التوصيل} \textsuperscript{(21)} بتاعته المشحونة \VUgras{براميل بيرة}
\textsuperscript{(22)} قدام فرشتها بالظبط، بيمنع زبوناتها إنهم يقرّبوا. وواحدة
بتزرع خضار بتجيب لها مقطف صغير \VUgras{خرشوف} \textsuperscript{(23)}.

%%K t15-tit-2 sub
\VUsaut{4.0mm}
\VUpk{10.2pt}{40}{البيع والشرا.}
\VUsaut{3.0mm}

%%K t15-04-1 p
4. --- وفي السوق الحركة كبيرة، علشان جت ساعة \VUgras{المزاد} \textsuperscript{(24)}.
و\VUgras{ستّات البيوت} \textsuperscript{(26)}، اللي جايين هنا يشتروا، بيقفوا وهما
بيعدّوا قدام دكّان \VUgras{بتاعة الكرشة} \textsuperscript{(25)}، علشان يشتروا
\VUgras{كرشة} أو \VUgras{راس بتلّو}.

%%K t15-05-1 p
5. --- و\VUgras{شيف} \textsuperscript{(27)} تخين بيفاصل على \VUgras{تونة} حلوة أوي
عند \VUgras{بتاعة السمك} \textsuperscript{(28)}، اللي بترفع أسعارها على كيفها. وهي
وصلها كمية كبيرة من \VUgras{الحنكليس}، و\VUgras{سمك موسى}، و\VUgras{التروتة}،
و\VUgras{المبروك}. و\VUgras{البكلاه} و\VUgras{بلح البحر} بتوعها طازة أوي.

%%K t15-tit-3 sub
\VUsaut{4.0mm}
\VUpk{10.2pt}{40}{بضاعة رخيصة وبضاعة غالية.}
\VUsaut{3.0mm}

%%K t15-06-1 p
6. --- و\VUgras{بتاعة الفراخ} \textsuperscript{(29)}، اللي فارشة جنبها، عندها ضمير
أوي. ولما بتبيع \VUgras{ديك رومي}، أو \VUgras{وزّة}، أو \VUgras{بطّة}، أو حتى
\VUgras{فرخة} عادية، ما بتطلبش غير السعر العادل.

%%K t15-07-1 p
7. --- وفي ناصية الميدان، في الدور الأول، من مدة قريّبة فتح \VUgras{بتاع قماش}
\textsuperscript{(30)}، اللي الزبونات عنده أكيد هيلاقوا تشكيلة حلوة أوي من
\VUgras{المفارش}، و\VUgras{الفوط}، و\VUgras{المراييل}، و\VUgras{الخِرَق}. وفي نفس
الدور، \VUgras{مسنّاتي} \textsuperscript{(31)} فتح ورشته. وهو بيسنّ \VUgras{سكاكين}
البايعات في السوق وبيعمل للجناينية \VUgras{مناجل} من أقوى \VUgras{صلب}.

%%K t15-tit-4 sub
\VUsaut{4.0mm}
\VUpk{10.2pt}{40}{البقالة.}
\VUsaut{3.0mm}

%%K t15-08-1 p
8. --- وتحت ورشته، من مدة قريّبة، فتح محل أكل كبير. وهو ما بقاش
\VUgras{البقّالة} \textsuperscript{(32)} البسيطة اللي ما لهاش أهمية زمان، اللي ما كانت
بتبيع غير البضاعة الشايعة، \VUgras{الشاي} \textsuperscript{(33)}،
و\VUgras{الشيكولاتة} \textsuperscript{(34)}، و\VUgras{السكر القوالب} \textsuperscript{(37)}،
و\VUgras{الصابون} \textsuperscript{(38)}. هنا بيبيعوا أي بضاعة، \VUgras{بسكوت}،
و\VUgras{معلّبات} \textsuperscript{(35)}، و\VUgras{مشروبات} \textsuperscript{(36)}،
و\VUgras{روم}، و\VUgras{كونياك}، و\VUgras{كرش}، و\VUgras{يانسون}،
و\VUgras{كوراساو}. وهناك بتلاقي طرايد: \VUgras{أرنب برّي} \textsuperscript{(39)}،
و\VUgras{سمنة} \textsuperscript{(40)}، و\VUgras{حجل} \textsuperscript{(41)}،
و\VUgras{سمّان} \textsuperscript{(42)}، و\VUgras{بطّ برّي} \textsuperscript{(43)}، أو
\VUgras{تدرّج} \textsuperscript{(44)}، بنفس السهولة اللي بتلاقي بيها فاكهة غريبة زي
\VUgras{الموز} \textsuperscript{(46)} و\VUgras{الأناناس}.

%%K t15-09-1 p
9. --- و\VUgras{صبي} \textsuperscript{(45)} البقّال، اللي بيحبّ يخدم أوي، مستعدّ
يديك اللي إنت عايزه: \VUgras{مربّى} \textsuperscript{(47)} بالكيلو أو بالربع،
و\VUgras{سمنة} \textsuperscript{(48)}، و\VUgras{خيار مخلّل} \textsuperscript{(49)}، أو
\VUgras{سردين} \textsuperscript{(50)} مملّح.

%%K t15-09-2 p
وده مريح أوي لـ\VUgras{الزبونات} \textsuperscript{(51)}، اللي بيلاقوا، جنب البضاعة
الشايعة، أندر الحاجات المبكّرة وأحلى الفاكهة: \VUgras{كمّترى} \textsuperscript{(52)}
مليان عصير، و\VUgras{برقوق} \textsuperscript{(53)}، و\VUgras{عنب} \textsuperscript{(54)}،
و\VUgras{خوخ} \textsuperscript{(55)}، و\VUgras{لوز} \textsuperscript{(56)}، و\VUgras{عين
جمل} \textsuperscript{(57)}.

%%K t15-tit-5 sub
\VUsaut{4.0mm}
\VUpk{10.2pt}{40}{الزباين.}
\VUsaut{3.0mm}

%%K t15-10-1 p
10. --- و\VUgras{الرز} \textsuperscript{(58)} و\VUgras{العدس} \textsuperscript{(59)} جنب
صندوق \VUgras{الرنجة} \textsuperscript{(60)} المدخّنة؛ و\VUgras{البطاطس}
\textsuperscript{(61)} و\VUgras{الفاصوليا} \textsuperscript{(63)} جنب \VUgras{التفاح}
\textsuperscript{(62)}، و\VUgras{التين} \textsuperscript{(64)}، و\VUgras{القراصيا}
\textsuperscript{(65)}، و\VUgras{البندق} \textsuperscript{(66)}. وبتلاقي في المحل ده
\VUgras{بيض} \textsuperscript{(67)} طازة و\VUgras{زيتون} \textsuperscript{(68)}؛ علشان
كده \VUgras{المقاطف} \textsuperscript{(70)} و\VUgras{شنط الشبك} \textsuperscript{(73)}
بتتملي بسرعة رهيبة.

%%K t15-11-1 p
11. --- و\VUgras{البنّ} \textsuperscript{(72)} بتاع المحل الممتاز ده كسب، بين
\VUgras{الشغّالات} \textsuperscript{(69)} و\VUgras{الطبّاخات}، سمعة يستاهلها، علشان
\VUgras{البقّال} \textsuperscript{(71)} العجوز نفسه بيحمّصه بعناية شديدة.

%%K t15-tit-6 sub
\VUsaut{4.0mm}
\VUpk{10.2pt}{40}{المناظر.}
\VUsaut{3.0mm}

%%K t15-12-1 p
12. --- ومش كل الناس بتيجي السوق علشان تتموّن. وفي الحركة الحلوة زي الصورة على
الحارة اللي بتودّي للشادر، بتشوف ناس من كل نوع. جنب \VUgras{صبي جزّار}
\textsuperscript{(75)} لطيف، بيدفع قدامه \VUgras{عربية إيد} \textsuperscript{(74)}، أهو
اتنين عساكر في أجازة فخورين أوي إنهم يورّوا بدلتهم اللمّاعة:
\VUgras{الهوسار} \textsuperscript{(76)} بـ\VUgras{دولمانه} الأزرق و\VUgras{شاكوه
المريّش}، و\VUgras{عريف الزواف}، بكمّه المنفوخ وراسه المغطّاة بـ\VUgras{طربوش}.

%%K t15-13-1 p
13. --- و\VUgras{الفرّان} \textsuperscript{(80)}، الواقف على عتبة \VUgras{الفرن}
\textsuperscript{(78)}، لسه عاجن عجينة \VUgras{عيشه} \textsuperscript{(79)} وطلّع من الفرن
\VUgras{الكرواسون} \textsuperscript{(81)}، و\VUgras{العيش الفينو} \textsuperscript{(82)}،
و\VUgras{الأرغفة} \textsuperscript{(83)} اللي في الفاترينة.

%%K t15-14-1 p
14. --- وفوق الفرن ساكنة \VUgras{خيّاطة} \textsuperscript{(84)} شاطرة، بتشغّل كام
\VUgras{مطرّزة} \textsuperscript{(85)}. وجنبها ساكنة \VUgras{مكوجيّة}
\textsuperscript{(86)}، اللي بتنشّي \VUgras{الغسيل} \textsuperscript{(88)} بعد ما تغسله
وتنشّفه، وبتكويه بـ\VUgras{المكواة} \textsuperscript{(87)}.

%%K t15-tit-7 sub
\VUsaut{4.0mm}
\VUpk{10.2pt}{40}{لحمة وسمك وخضار.}
\VUsaut{3.0mm}

%%K t15-15-1 p
15. --- وفي \VUgras{الجزارة} \textsuperscript{(89)}، اللي في الدور الأرضي، بيبيعوا كل
أنواع \VUgras{اللحمة}: \VUgras{لحمة ضاني} \textsuperscript{(90)}، أو \VUgras{لحمة
بقري} \textsuperscript{(94)}، أو \VUgras{لحمة بتلّو} \textsuperscript{(92)}، على شكل
\VUgras{فخاد ضاني، وستيك، وروستو،} و\VUgras{كستليتة}. و\VUgras{صبي الجزّار}
\textsuperscript{(93)} بيقطّع اللحمة دي كلها وبيحطّ لوحده \VUgras{عضم النخاع}
\textsuperscript{(96)}. وعلى باب الجزارة في \VUgras{بتاعة الجندوفلي}
\textsuperscript{(97)} بتبيع \VUgras{جندوفلي} \textsuperscript{(98)}. وبتفتحه لك لو
إنت عايز.

%%K t15-16-1 p
16. --- والتجّار دول كلهم بيدفعوا رخصة أو رسم مكان؛ علشان كده \VUgras{العسكري}
\textsuperscript{(100)} بيطارد من غير رحمة \VUgras{بتاعات الخضار المتجوّلات}
\textsuperscript{(99)} المساكين اللي بينافسوهم وهما بيوزّعوا على عربياتهم
\VUgras{جزر، وفجل، ولفت، وكرّات،} أو \VUgras{حزم هليون، وبسلّة طازة، وسبانخ،
وحمّيض، وجرجير} و\VUgras{بقدونس}.

%%K t15-tit-8 sub
\VUsaut{4.0mm}
\VUpk{10.2pt}{40}{خد بالك من العسكري!}
\VUsaut{3.0mm}

%%K t15-17-1 p
17. --- والولد اللي بيبيع \VUgras{توم} \textsuperscript{(101)} و\VUgras{بصل}، عارف
كويس أوي إنه هو كمان مش مسموح له

%%K t15-17-1 p suite
يبيع بضاعته جنب السوق. وهو بيهرب، وبيجري وهو معرّض إنه يقلب مقطف
\VUgras{الكابوريا}، و\VUgras{السلطعون}، و\VUgras{الجمبري} بتاع \VUgras{بتاع}
\textsuperscript{(102)} \VUgras{الاستاكوزا} و\VUgras{الأمّ جراب}.

%%K t15-18-1 p
18. --- والكل بيتحرّك وبيعمل دوشة؛ بس ده ما بيمنعش إنك تسمع بوضوح صوت
\VUgras{القزّاز} \textsuperscript{(103)} المتجوّل، ولا نداء \VUgras{الفحّام}
\textsuperscript{(104)} اللي ساب عربيته كام لحظة علشان يوصّل \VUgras{شكارة فحم}
\textsuperscript{(105)}.

\VUsaut{6.0mm}
\VUfilet{22mm}
